\subsection{Diagonal arithmetic points and weight-zero rigidity}
\label{subsec:selmer-diagonal-rigidity}

Return to the arithmetic setting. Let $K$ be an imaginary quadratic
field, let $\ell$ be an odd prime, and let
$E=\overline{\mathbf Q}_\ell$, with valuation ring $\mathcal O_E$
and maximal ideal $\mathfrak m_E$. Write
\[
 \Psi=\chi_\ell^e\psi,\qquad e\ge1,\qquad
 (w_1,w_2,w_3)=(1,0,-1),
\]
where $\psi:G_{\mathbf Q}\to E^\times$ has finite order prime to
$\ell$. The prescribed representation on $G_K$ is
$\rho_0=\operatorname{diag}(\Psi,1,\Psi^{-1})$.

\begin{lemma}[Integral comparison on finite groups]
\label{lem:selmer-inertia-averaging}
Let $R$ be a commutative local ring and let $H$ be a finite group
whose order is a unit in $R$. If
$\rho,\rho_0:H\to\mathrm{GL}_n(R)$ have the same residual
representation, then
\[
 M=\frac1{|H|}\sum_{h\in H}\rho_0(h)\rho(h)^{-1}
\]
is invertible, reduces to the identity, and satisfies
$\rho_0(g)M=M\rho(g)$. It commutes with coefficient homomorphisms.
If both representations are diagonal, they are equal.
\end{lemma}

\begin{proof}
Every summand reduces to the identity, so $\det M$ is a unit.
Replacing $h$ by $gh$ in
$\rho_0(g)M\rho(g)^{-1}$ proves the intertwining identity.
The finite formula commutes with coefficient maps. In the diagonal
case, $M$ is diagonal with unit entries, and the identity gives
coordinatewise equality.
\end{proof}

\begin{proposition}[Uniqueness of a diagonal arithmetic point]
\label{prop:selmer-diagonal-point}
Assume $\ell\nmid h_K$. Let
$\rho|_{G_K}=\operatorname{diag}(\alpha_1,\alpha_2,\alpha_3)$
be a continuous representation with unit entries, residually equal
to $\rho_0$. Suppose that:
\begin{enumerate}
\item for each finite $v\nmid\ell$, $\rho$ kills
      $\ker(\psi|_{I_v})$;
\item for each $v\mid\ell$,
      $\operatorname{tr}\rho|_{I_v}=\Psi+1+\Psi^{-1}$, and
      some $s_v\in I_v$ satisfies
      $\overline{\Psi(s_v)}^{\,2}\ne1$.
\end{enumerate}
Then $\alpha_i=\Psi^{w_i}$ for every $i$.
If the two polarized representations of $G_{\mathbf Q}$ also
have the same prescribed value at complex conjugation, they are equal.
\end{proposition}

\begin{proof}
For $v\nmid\ell$, the cyclotomic character is unramified, so
both inertia representations factor through the same finite image
of $\psi$. Its order is prime to $\ell$.
Lemma~\ref{lem:selmer-inertia-averaging} proves coordinatewise
equality on $I_v$.

For $v\mid\ell$, distinct characters of $I_v$ are linearly
independent: a shortest nonzero relation, translated by an element
distinguishing two characters and subtracted from a scalar multiple
of itself, would give a shorter nonzero relation.
The trace identity therefore identifies the multisets
$\{\alpha_1,\alpha_2,\alpha_3\}$ and $\{\Psi,1,\Psi^{-1}\}$
on $I_v$. Characteristic zero ensures equality of their
multiplicities. At $s_v$, the three residual values are pairwise
distinct, so the residual identities force the ordered matching
$\alpha_i|_{I_v}=\Psi^{w_i}|_{I_v}$.

Thus $\beta_i=\alpha_i\Psi^{-w_i}$ is everywhere unramified
and takes values in $1+\mathfrak m_E$.
Hilbert class field reciprocity identifies the maximal unramified
abelian quotient of $G_K$ with its finite ideal class group
\cite{milne-cft}. Hence $\beta_i^{h_K}=1$.
This also holds for $E$-valued characters: their compact,
totally disconnected images are profinite, and every finite
quotient factors through the same class group.
If $u\in1+\mathfrak m_E$ and $u^{h_K}=1$, then
\[
 (u-1)(1+u+\cdots+u^{h_K-1})=0.
\]
The second factor is a unit, since it reduces to $h_K$.
Therefore $\beta_i=1$.
Finally $G_{\mathbf Q}=G_K\sqcup G_Kc$ for complex
conjugation $c$, so agreement on $G_K$ and at $c$ suffices.
\end{proof}

For the prescribed family, $\ell\nmid\operatorname{disc}K$,
$\psi^2|_{I_v}=1$ at $v\mid\ell$, and
$\ell-1\nmid2e$ give the separating element directly.
Indeed, inertia has a surjective mod-$\ell$ cyclotomic character
and $\overline\Psi^{\,2}=\overline\chi_\ell^{\,2e}$ there.
An ordinary invariant flag with ordered inertia quotients
$\Psi,1,\Psi^{-1}$ supplies the required trace.
This local ordinary statement remains an input; prescribed
Hodge--Tate weights alone are not being used to infer it.
For a character $\psi$ of order a power of $2$, its order is
automatically prime to odd $\ell$; the other exclusions remain.

The proposition gives the fixed-point separation of the graded model
once its point comparison is supplied. More precisely, retain the
lifting of each equation point, up to the torus action, to a local
point of the marked representation ring satisfying the proposition's
residual and local conditions; the identification of
negative and positive vanishing with upper and lower triangularity;
and the determination of a model point by its representation.
A point vanishing on every nonzero weight then lifts to a diagonal
representation, which the proposition identifies with the base
representation. Its model point is the augmentation, fixed by the
torus. The lifting and determination properties, including their
finiteness and density hypotheses, remain separate requirements.

\begin{proposition}[Weight-zero cocycles]
\label{prop:selmer-zero-weight-rigidity}
Let $V$ be a finite-dimensional continuous $E$-representation of
$G_{\mathbf Q}$ on which $G_K$ acts trivially.
If a continuous cocycle $\gamma:G_{\mathbf Q}\to V$ vanishes
on every finite inertia group of $K$, then it is a coboundary,
with primitive
\[
 m=-\frac12\gamma(c),\qquad \gamma(g)=gm-m.
\]
If $\gamma(c)=0$, then $\gamma=0$. This statement does not
require $\ell\nmid h_K$.
\end{proposition}

\begin{proof}
On $G_K$ the cocycle is an additive homomorphism. Its coordinates
factor through the finite class group; the additive group of $E$
has no nonzero finite subgroup, so the restriction vanishes.
On the other coset its value is $\gamma(c)$.
The identity $c^2=1$ gives $c\gamma(c)=-\gamma(c)$, and hence
$cm-m=\gamma(c)$ for the stated $m$.
\end{proof}

The weight-zero adjoint space of $\rho_0$ consists of diagonal
matrices, on which $G_K$ acts trivially. The last proposition
therefore proves its rigidity once the actual local finite Selmer
condition is identified with unramified cohomology for the trivial
character. The passage from ring derivations to those Selmer
cocycles is also retained. Neither characteristic-zero point
uniqueness nor this tangent statement removes integral conormal
torsion. The next subsection supplies negative-weight base rigidity.
Rigidity at two-sided points and construction of the marked arithmetic
model remain separate requirements.
